\begin{table}[htbp]
\centering
\caption{Discipline support in the balanced analysis cohort. Codes 8--11 (relay, throws, race walk, and one further unused code) had zero respondents and were excluded from discipline-level adjustment.}
\label{tab:discipline-support}
\begin{tabular}{llrl}
\toprule
Code & Discipline & n & Included in model \\
\midrule
1 & Sprint & 107 & Yes \\
2 & Middle distance & 69 & Yes \\
3 & Long distance & 59 & Yes \\
4 & Hurdles & 96 & Yes \\
5 & Horizontal jumps & 44 & Yes \\
6 & Pole vault & 19 & Yes \\
7 & Combined events & 41 & Yes \\
8--11 & (zero respondents) & 0 & No \\
\bottomrule
\end{tabular}
\end{table}

\begin{table}[htbp]
\centering
\caption{FIDAL competition-registry counts used to construct the 2026 eligible-population frame. Under 20 counts for 2026 were used directly. Because 2026 counts were not yet available for Under 23 and Senior athletes when recruitment was planned, their sex-specific frame estimates were the rounded means of the 2023--2025 counts. The resulting age--sex frame totaled 1328 athletes; 2025 discipline shares within each age--sex cell were used to allocate stratum targets.}
\label{tab:population-frame-history}
\begin{tabular}{llrrrrr}
\toprule
Age group & Sex & 2023 & 2024 & 2025 & 2026 & Frame estimate \\
\midrule
Under 20 & Women & 235 & 265 & 274 & 204 & 204 \\
Under 20 & Men & 327 & 361 & 338 & 311 & 311 \\
Under 23 & Women & 217 & 211 & 207 & -- & 212 \\
Under 23 & Men & 244 & 247 & 252 & -- & 248 \\
Senior & Women & 176 & 176 & 156 & -- & 169 \\
Senior & Men & 195 & 199 & 159 & -- & 184 \\
\midrule
Total & & 1394 & 1459 & 1386 & -- & 1328 \\
\bottomrule
\end{tabular}
\end{table}

\begin{table}[htbp]
\centering
\caption{Bayesian worst-case population-proportion precision calculation used to set Phase 2 (targeted recall) enrollment targets. The calculation set a minimum recruitment target; representativeness was addressed separately through stratified recall and poststratification.}
\label{tab:sample-size}
\begin{tabular}{lr}
\toprule
Quantity & Value \\
\midrule
Planning quantity & Population proportion \\
Prior & Beta(1,1) \\
Worst-case count & $x=\lfloor n/2\rfloor$ \\
Posterior at candidate $n$ & Beta$(1+x,1+n-x)$ \\
Interval and target & Equal-tailed 95\% CrI; half-width $\leq 0.06$ \\
$n=263$ half-width & 0.060038 \\
$n=264$ half-width & 0.059926 \\
Estimated eligible population & 1328 \\
Minimum target sample & 264 (19.9\% of frame) \\
Consented, eligible, complete & 355 (26.7\% of frame) \\
Balanced analysis cohort & 352 (26.5\% of frame) \\
Target allocation & Proportional across 42 strata \\
\bottomrule
\end{tabular}
\end{table}

\begin{table}[htbp]
\centering
\caption{Age--sex poststratification weight diagnostics (Section~\ref{sec:weighting}). Model weight is the regularized ratio of target population share to sample share; Kish effective sample size after weighting was 340.5 (97\% of $N=352$).}
\label{tab:weight-diagnostics}
\begin{tabular}{llrrrr}
\toprule
Age group & Sex & Target $n$ & Sample $n$ & Sample share & Model weight \\
\midrule
Under 20 & Women & 204 & 43 & 12.2\% & 1.26 \\
Under 23 & Women & 212 & 53 & 15.1\% & 1.06 \\
Senior & Women & 169 & 62 & 17.6\% & 0.72 \\
Under 20 & Men & 311 & 71 & 20.2\% & 1.16 \\
Under 23 & Men & 248 & 63 & 17.9\% & 1.04 \\
Senior & Men & 184 & 60 & 17.0\% & 0.81 \\
\bottomrule
\end{tabular}
\end{table}

\begin{table}[htbp]
\centering
\caption{Age- and sex-conditional posterior estimate of the competition-period dose association (75th vs.\ 25th percentile carbon share among users) from the primary HSGP model, holding the model structure and covariate adjustment fixed. $P_{\mathrm{dir}}$ is the posterior probability of the median estimate's direction. These are descriptive, hypothesis-generating comparisons (Section~\ref{sec:subgroups}), not independently powered tests of effect modification.}
\label{tab:subgroups}
\begin{tabular}{llrr}
\toprule
Grouping & Level & Difference, pp (95\% CrI) & $P_{\mathrm{dir}}$ \\
\midrule
Sex & Men & 14.0 ($-0.2$ to 30.6) & 96.8\% \\
Sex & Women & 13.1 ($-0.1$ to 29.1) & 96.8\% \\
Age group & Under 20 & 13.1 ($-0.1$ to 28.8) & 96.8\% \\
Age group & Under 23 & 14.0 ($-0.2$ to 30.6) & 96.8\% \\
Age group & Senior & 14.1 ($-0.2$ to 30.8) & 96.8\% \\
\bottomrule
\end{tabular}
\end{table}

\begin{table}[htbp]
\centering
\caption{Population-weighted versus sample (unweighted) adjusted estimands from the primary HSGP model. Weighting applies the age--sex poststratification of Table~\ref{tab:weight-diagnostics}. $P_{\mathrm{dir}}$ is reported beside every 95\% CrI.}
\label{tab:weighted-vs-unweighted}
\small
\begin{tabular}{llrrrr}
\toprule
Period & Contrast & Unweighted, pp (95\% CrI) & $P_{\mathrm{dir}}$ & Weighted, pp (95\% CrI) & $P_{\mathrm{dir}}$ \\
\midrule
Preparation & Adoption difference & 6.5 ($-5.6$ to $19.2$) & 85.3\% & 4.8 ($-7.7$ to $16.8$) & 77.6\% \\
Preparation & Dose difference & 0.0 ($-12.0$ to $10.2$) & 50.6\% & 0.3 ($-10.4$ to $10.4$) & 54.5\% \\
Competition & Adoption difference & 4.7 ($-9.6$ to $18.6$) & 74.7\% & 7.1 ($-7.4$ to $20.6$) & 84.0\% \\
Competition & Dose difference & 13.2 ($-0.2$ to 29.1) & 96.9\% & 13.7 ($-0.2$ to 30.0) & 96.8\% \\
\bottomrule
\end{tabular}
\end{table}

\begin{table}[htbp]
\centering
\caption{Model comparison and validation diagnostics for the primary HSGP (exposure) and matched baseline (no-exposure) models. ELPD: expected log predictive density (athlete-grouped PSIS-LOO); all Pareto-$k \leq 0.5$ for both models.}
\label{tab:model-comparison}
\begin{tabular}{lrrrr}
\toprule
Model & ELPD$_{\text{LOO}}$ & SE & Max $\hat{R}$ & Divergences \\
\midrule
Primary (exposure) & $-429.00$ & 11.14 & 1.005 & 0 \\
Baseline (no exposure) & $-430.92$ & 11.03 & 1.007 & 0 \\
\midrule
\multicolumn{5}{l}{ELPD difference (exposure $-$ baseline): 1.92 (SE 2.41)} \\
\bottomrule
\end{tabular}
\end{table}
